%! AP Statistics — Unit 1, Lesson 1.5 — Homework ANSWER KEY
\documentclass[10pt]{article}
\usepackage{apstats-article}
\usepackage{apstats-key}

%=============================================================================
\begin{document}

\pageheader{Unit 1, Lesson 1.5}{Homework --- Answer Key}
\begin{tabularx}{\linewidth}{X X X}
Name:\;\ans{Key} & Date:\; & Period:\;
\end{tabularx}

\vspace{0.1in}

% ── PART A KEY ───────────────────────────────────────────────────────────────
\begin{blurbbox}[Part A \quad Interpret a Dotplot --- Key]{navylight}
\small

\begin{center}
\begin{tikzpicture}[scale=1.05]
  \draw[thick] (4.2,0) -- (11.8,0);
  \foreach \x/\lbl in {5/5, 6/6, 7/7, 8/8, 9/9, 10/10, 11/11} {
    \draw (\x,-0.12) -- (\x,0.12);
    \node[below, font=\small] at (\x,-0.12) {\lbl};
  }
  \node[below, font=\small\itshape] at (8,-0.52) {Time (minutes)};
  \filldraw[navy] (5, 0.38) circle (0.14);
  \foreach \k in {1,...,3} \filldraw[navy] (6, \k*0.38) circle (0.14);
  \foreach \k in {1,...,5} \filldraw[navy] (7, \k*0.38) circle (0.14);
  \foreach \k in {1,...,5} \filldraw[navy] (8, \k*0.38) circle (0.14);
  \foreach \k in {1,...,3} \filldraw[navy] (9, \k*0.38) circle (0.14);
  \foreach \k in {1,...,2} \filldraw[navy] (10, \k*0.38) circle (0.14);
  \filldraw[navy] (11, 0.38) circle (0.14);
\end{tikzpicture}
\end{center}

\begin{enumerate}[leftmargin=1.5em, itemsep=6pt]

  \item \ans{Range $= 11 - 5 = 6$ minutes.}
        \ans{Students finishing in $\le 8$ min: $1+3+5+5 = 14$ students.}

  \item \ansline{The distribution is approximately symmetric. Both sides of the
        peak (7 and 8 minutes) have matching counts: 5 students at 7 min, 5 at 8 min,
        with 3 on each side (6 and 9), then 1--2 tapering off at the extremes.}

  \item \ans{Students taking more than 9 min: $2 + 1 = 3$ students.
        $3/20 = 0.15 = 15\%$.}

  \item \ansline{The distribution of problem-set completion times for these 20
        students is approximately symmetric, centered around 7--8 minutes (the
        most common times, each with 5 students). Times range from 5 to 11 minutes,
        a spread of 6 minutes, with most students finishing between 6 and 9 minutes.}

  \begin{teachernote}
  Full credit requires: shape (symmetric), center (7 or 8 minutes), spread (range
  or interval), and context (students, problem set). Deduct for missing context
  or citing raw count instead of a center measure.
  \end{teachernote}

\end{enumerate}
\end{blurbbox}

\vspace{0.1in}

% ── PART B KEY ───────────────────────────────────────────────────────────────
\begin{blurbbox}[Part B \quad Stemplot --- Key]{greenacc}
\small

\begin{center}
\renewcommand{\arraystretch}{1.4}
\begin{tabular}{r | l}
\textbf{Stem} & \textbf{Leaf} \\
\hline
5 & \ans{2\; 5\; 8} \\
6 & \ans{1\; 3\; 5\; 7\; 8} \\
7 & \ans{0\; 2\; 3\; 4\; 6\; 8} \\
8 & \ans{2} \\
\end{tabular}
\end{center}

\begin{enumerate}[leftmargin=1.5em, itemsep=6pt]

  \item \ans{$n = 15$, so the median is the 8th value. Reading the stemplot in order:
        52, 55, 58, 61, 63, 65, 67, 68, 70, 72, 73, 74, 76, 78, 82.
        The 8th value is \textbf{68$^\circ$F}.}

  \item \ansline{The distribution is roughly symmetric or slightly skewed right.
        Most days had temperatures in the 60s--70s, with the single value of 82$^\circ$F
        pulling the right tail slightly. This suggests October temperatures were
        mild with occasional warmer days.}

  \item \ans{Sorted values: 52, 55, 58, 61, 63, \underline{65}, 67, 68, 70, \underline{72}, 73, 74, 76, 78, 82.
        $n=15$: Q1 is median of lower half (positions 1--7): 4th value = 61$^\circ$F.
        Q3 is median of upper half (positions 9--15): 12th value = 74$^\circ$F.
        IQR $= 74 - 61 = \mathbf{13}^\circ$F.}

  \begin{teachernote}
  For $n=15$: Q1 = median of {52,55,58,61,63,65,67} = 61;
  Q3 = median of {70,72,73,74,76,78,82} = 74. IQR = 13.
  Accept answers using either the exclusive or inclusive quartile method if
  work is shown consistently.
  \end{teachernote}

\end{enumerate}
\end{blurbbox}

\vspace{0.1in}

% ── PART C KEY ───────────────────────────────────────────────────────────────
\begin{blurbbox}[Part C \quad AP-Style --- Key]{redacc}
\small

\renewcommand{\arraystretch}{1.6}
\begin{tabularx}{0.75\linewidth}{@{} C{0.26\linewidth} C{0.22\linewidth} C{0.26\linewidth} @{}}
\rowcolor{sky}
\textbf{Time (min)} & \textbf{Frequency} & \textbf{Rel.\ Frequency} \\
\midrule
$[0, 10)$   & 2  & \ans{0.050} \\
$[10, 20)$  & 8  & \ans{0.200} \\
$[20, 30)$  & 14 & \ans{0.350} \\
$[30, 40)$  & 11 & \ans{0.275} \\
$[40, 50)$  & 5  & \ans{0.125} \\
\midrule
\rowcolor{sky}\textbf{Total} & 40 & \ans{1.000} \\
\end{tabularx}

\vspace{0.1in}

\begin{enumerate}[leftmargin=1.5em, itemsep=6pt]

  \item (Answers in table above.)

  \item \ans{The $[20, 30)$ bin contains the most students: 14 students, or 35\%.}

  \item \ans{$(11 + 5)/40 = 16/40 = 0.40 = 40\%$ spent at least 30 minutes.}

  \item \ansline{The distribution of homework time for these 40 students is
        unimodal and skewed right, with a peak in the 20--30 minute bin (35\% of
        students). A typical student spent roughly 20--30 minutes on homework, and
        the times range from under 10 minutes to nearly 50 minutes (a spread of
        about 40 minutes), with the right tail indicating a few students who spent
        substantially more time than the rest.}

  \begin{teachernote}
  Require: shape (skewed right or unimodal), a measure of center stated in context,
  a measure of spread with units, and use of context (students, homework time).
  Deduct if the student does not mention the direction of skew or gives a center
  without units.
  \end{teachernote}

  \item \ansline{No, a stemplot would not be better here. With $n = 40$ data points
        and values ranging from 0 to nearly 50 minutes, a stemplot would have many
        leaves on each stem and could become difficult to read. A histogram
        efficiently shows the overall shape of the distribution for a dataset of
        this size. A stemplot is best for $n \le 30$--50 when preserving individual
        values is important.}

\end{enumerate}
\end{blurbbox}

\end{document}
